\chapter{Molekul, Obat, dan Data Medis}
\label{ch:hayati}

Kuliah AI in life sciences saya ikuti pada musim panas 2025. Materinya luas: penemuan obat dan
representasi molekul, QSAR dan virtual screening, geometric deep learning untuk molekul, model
multimodal, model generatif, data medis, sampai perencanaan sintesis kimia. Banyak metode yang
dibahas lahir di Linz sendiri, dari DeepTox sampai jaringan Hopfield modern untuk kimia, sehingga
kuliahnya banyak memperlihatkan hasil riset dari kelompok-kelompok di kampus yang sama.

Menemukan obat baru memakan waktu lebih dari satu dekade dan biaya yang amat besar, dan sebagian
besar kandidat gagal di tengah jalan. Pipeline-nya panjang: memilih \istilah{target} biologis, biasanya
sebuah protein yang terlibat dalam penyakit, mencari \istilah{hit} yang memengaruhi target itu,
mengembangkannya menjadi \istilah{lead}, mengoptimalkan sifat-sifatnya, lalu menguji pada hewan dan
manusia. Machine learning tidak bisa menggantikan uji klinis, tetapi bisa membuat langkah-langkah awal
jauh lebih murah, misalnya dengan menyaring jutaan molekul di komputer sebelum satu pun disintesis.

\section{Bagaimana mesin melihat molekul}

Molekul organik kecil adalah atom-atom yang dihubungkan ikatan kimia. Representasi paling alami adalah
graf: atom sebagai simpul, ikatan sebagai sisi, masing-masing dengan atribut seperti jenis unsur,
muatan, dan orde ikatan. Untuk disimpan dan dicari di basis data, graf ini ditulis sebagai string.
SMILES \citep{weininger1988} berjalan menyusuri graf dan mencatat atom yang dilewati: etanol ditulis
\texttt{CCO}, benzena \texttt{c1ccccc1}, dengan angka yang menandai titik pembukaan dan penutupan
cincin. Satu molekul bisa punya banyak SMILES yang sah, tergantung dari atom mana penyusuran dimulai,
sehingga dibutuhkan bentuk kanonik untuk membandingkan. SELFIES dirancang agar setiap string yang mungkin
selalu mewakili molekul yang sah, sifat yang berguna untuk model generatif. InChI dan InChIKey memberi
pengenal unik yang bisa dipakai sebagai kunci basis data.

Untuk machine learning klasik, molekul diubah menjadi vektor fitur. \istilah{Molecular descriptor} adalah
besaran yang bisa dihitung, seperti berat molekul, jumlah donor ikatan hidrogen, atau koefisien partisi
oktanol--air ($\log P$) yang mengukur seberapa ``berminyak'' sebuah molekul. \istilah{Molecular
fingerprint} adalah vektor biner panjang yang menandai keberadaan pola-pola substruktur. Extended
connectivity fingerprint \citep{rogers2010} dibangun dengan cara yang mirip uji Weisfeiler--Leman
di \cref{ch:gdl}. Setiap atom mulai dengan pengenal dari sifat-sifatnya sendiri. Di setiap iterasi,
pengenal atom diperbarui dengan meng-\emph{hash} pengenalnya bersama pengenal tetangga-tetangganya,
sehingga setelah dua atau tiga iterasi, setiap pengenal mewakili lingkungan berjari-jari dua atau tiga
ikatan. Semua pengenal kemudian dilipat ke dalam vektor berukuran tetap, misalnya 2048 bit. ECFP
adalah message passing tanpa parameter yang dipelajari.

Kemiripan dua molekul diukur dengan koefisien Tanimoto pada fingerprint-nya. Koefisien ini adalah
rasio irisan terhadap gabungan,
\begin{equation}
T(A,B)=\frac{|A\cap B|}{|A\cup B|}=\frac{c}{a+b-c},
\end{equation}
dengan $a$ dan $b$ jumlah bit yang menyala pada masing-masing fingerprint dan $c$ jumlah bit yang
menyala pada keduanya. Bentuk kedua didapat karena $|A\cup B|=|A|+|B|-|A\cap B|$: bit yang ada di
keduanya terhitung dua kali dalam $a+b$, sehingga harus dikurangi sekali. Kalau molekul $A$ punya 20
bit menyala, $B$ punya 18, dan 12 di antaranya sama, kemiripannya $12/(20+18-12)=12/26\approx0{,}46$.
Dalam praktik, nilai di atas sekitar $0{,}7$ untuk ECFP sering dianggap tanda dua molekul cukup mirip
untuk kemungkinan punya aktivitas serupa, tetapi ambang ini bergantung pada jenis fingerprint.

\section{Aktivitas biologis}

Seberapa kuat sebuah senyawa memengaruhi targetnya diukur dalam percobaan yang disebut \istilah{assay}.
Pada konsentrasi yang semakin tinggi, efeknya biasanya naik mengikuti kurva berbentuk S. Bentuk ini
bisa diturunkan dari kesetimbangan pengikatan yang paling sederhana. Misalkan reseptor $R$ dan ligan
$L$ membentuk kompleks $RL$ dengan konstanta disosiasi
\begin{equation}
K_d=\frac{[R][L]}{[RL]}.
\end{equation}
Jumlah reseptor total tetap, $[R]_{\mathrm{tot}}=[R]+[RL]$. Fraksi reseptor yang terikat adalah
$\theta=[RL]/[R]_{\mathrm{tot}}$. Dari definisi $K_d$, $[R]=K_d[RL]/[L]$, sehingga
\begin{equation}
\theta=\frac{[RL]}{K_d[RL]/[L]+[RL]}=\frac{[L]}{K_d+[L]}.
\end{equation}
Ketika konsentrasi ligan sama dengan $K_d$, tepat separuh reseptor terikat. Pada skala logaritmik
konsentrasi, fungsi ini menjadi kurva S yang terkenal. Untuk penghambatan, konsentrasi yang memberi
separuh efek maksimum disebut IC$_{50}$, dan untuk pemodelan biasanya dipakai
$\mathrm{pIC}_{50}=-\log_{10}(\mathrm{IC}_{50}\ \text{dalam molar})$. Senyawa dengan IC$_{50}$ 100
nanomolar punya pIC$_{50}$ sebesar 7, dan setiap kenaikan satu satuan berarti senyawa sepuluh kali
lebih kuat. Skala logaritmik ini membuat target regresi lebih mendekati distribusi normal, sesuai
asumsi noise Gauss dari \cref{ch:data}.

Obat tidak cukup kuat, ia juga harus sampai ke tempatnya. Sifat \istilah{ADME}, yaitu penyerapan,
distribusi, metabolisme, dan ekskresi, menentukan hal ini. Aturan lima Lipinski \citep{lipinski1997}
adalah heuristik terkenal untuk obat yang diminum: berat molekul tidak lebih dari 500 dalton, $\log P$
tidak lebih dari 5, paling banyak 5 donor dan 10 akseptor ikatan hidrogen. Aturan ini dirumuskan dari
pengamatan pada obat-obat yang sudah berhasil, bukan dari hukum fisika, dan banyak pengecualiannya.

Data aktivitas dikumpulkan dalam basis data publik seperti ChEMBL \citep{gaulton2012}. Kalau disusun
sebagai matriks senyawa terhadap target, matriks ini jarang: sebagian besar pasangan tidak
pernah diukur. Datanya juga tidak seimbang, karena senyawa aktif jauh lebih sedikit dari yang tidak aktif,
dan berisik, karena laboratorium berbeda memberi angka berbeda untuk pengukuran yang sama.

\section{QSAR dan virtual screening}

\istilah{Quantitative structure--activity relationship} (QSAR) adalah nama lama untuk apa yang kini kita
sebut supervised learning pada molekul: memprediksi aktivitas dari struktur. Pipeline-nya mengikuti
\cref{ch:data}: pilih representasi, pilih model, pilih hyperparameter dengan validasi, lalu evaluasi.
\istilah{Virtual screening} memakai model seperti ini, atau sekadar pencarian kemiripan, untuk menyaring
pustaka berisi jutaan sampai miliaran senyawa dan memilih segelintir yang layak diuji di laboratorium.
Pendekatan berbasis ligan hanya memakai senyawa aktif yang sudah diketahui. Pendekatan berbasis struktur
memakai bentuk tiga dimensi protein target untuk menilai seberapa baik sebuah senyawa bisa menempel.

Kuliah memberi penekanan besar pada satu jebakan: cara membagi data. Ahli kimia bekerja dalam seri
senyawa, yaitu banyak varian kecil dari satu kerangka yang sama. Kalau data dibagi acak, varian dari seri
yang sama masuk ke data latih dan data uji sekaligus, dan model tampak hebat karena pada dasarnya
hanya mengenali saudara dekat. Pembagian berdasarkan kerangka molekul atau berdasarkan klaster
memaksa model diuji pada kimia yang benar-benar baru, dan skornya biasanya turun drastis. Bias ini
disebut \istilah{compound series bias}. Pada model multitask, kebocoran bisa terjadi lewat tugas lain yang
mengukur senyawa yang sama, dan pada prediksi sinergi dua obat, lewat pasangan obat yang sama dalam
urutan terbalik. Validasi silang bersarang dipakai agar pemilihan hyperparameter tidak ikut membocorkan
data uji.

DeepTox \citep{mayr2016}, dari kelompok di Linz, menunjukkan kekuatan jaringan saraf multitask untuk
memprediksi toksisitas. Banyak tugas toksisitas berbagi mekanisme, sehingga satu jaringan yang belajar
semuanya bersama-sama membangun representasi yang lebih baik daripada model terpisah untuk setiap tugas.
ChemProp \citep{yang2019} membawa message passing langsung ke graf molekul, sehingga fiturnya dipelajari
alih-alih dirancang. Untuk sifat yang bergantung pada bentuk tiga dimensi, jaringan ekuivarian seperti
EGNN \citep{satorras2021} menjamin bahwa prediksi tidak berubah kalau molekul diputar.

\section{Belajar dari banyak modalitas}

Data ilmu hayati datang dalam banyak bentuk: struktur molekul, deskripsi teks sebuah assay, gambar
mikroskop sel, urutan protein. Kuliah membahas cara menghubungkannya dengan \istilah{contrastive learning}.
Gagasan dasarnya berasal dari CLIP \citep{radford2021}, yang melatih encoder gambar dan encoder teks
agar pasangan yang benar berdekatan dan pasangan yang salah berjauhan di ruang embedding bersama.

Loss-nya bisa diturunkan sebagai masalah klasifikasi. Dalam satu batch ada $N$ pasangan. Untuk molekul
ke-$i$, kita hitung kemiripan kosinus $s_{ij}$ dengan setiap deskripsi $j$ dalam batch, lalu
memperlakukan pertanyaan ``deskripsi mana yang cocok dengan molekul ini?'' sebagai klasifikasi
$N$ kelas dengan softmax. Cross-entropy untuk jawaban yang benar memberi
\begin{equation}
\Loss_i=-\log\frac{\exp(s_{ii}/\tau)}{\sum_{j=1}^N\exp(s_{ij}/\tau)},
\end{equation}
dengan $\tau$ suhu yang mengatur ketajaman softmax. Loss ini dihitung ke dua arah, dari molekul ke teks
dan dari teks ke molekul, lalu dirata-ratakan. Sebagai contoh, misalkan batch berisi tiga pasangan,
$\tau=1$, dan untuk molekul pertama kemiripannya $(0{,}9;\ 0{,}2;\ 0{,}1)$. Bobot softmax untuk pasangan
yang benar adalah $e^{0{,}9}/(e^{0{,}9}+e^{0{,}2}+e^{0{,}1})\approx2{,}46/(2{,}46+1{,}22+1{,}11)\approx0{,}51$,
dan loss-nya $-\ln0{,}51\approx0{,}67$. Menurunkan suhu menjadi $0{,}1$ membuat bobot itu mendekati satu.

CLAMP \citep{seidl2023} memakai gagasan ini untuk menghubungkan molekul dengan deskripsi teks assay,
sehingga model bisa memprediksi aktivitas untuk assay baru hanya dari deskripsinya, tanpa satu pun
contoh berlabel. Ini adalah \istilah{zero-shot learning}. Ketika hanya ada segelintir contoh berlabel,
pendekatan \istilah{few-shot} memakai jaringan Hopfield modern (\cref{ch:lstm}) untuk membandingkan molekul
kueri dengan contoh-contoh yang ada dan dengan konteks dari molekul lain. \istilah{Active learning} memilih
molekul berikutnya yang paling informatif untuk diukur, biasanya yang prediksinya paling tidak pasti.

\section{Merancang molekul baru}

Model generatif membalik arah QSAR: bukan dari struktur ke sifat, melainkan dari sifat yang diinginkan
ke struktur. Pendekatan paling sederhana memperlakukan SMILES sebagai teks dan melatih LSTM untuk menebak
karakter berikutnya \citep{segler2018gen}. Setelah dilatih pada jutaan molekul, model ini menghasilkan
string baru yang sebagian besar adalah molekul sah. Untuk mengarahkannya ke sifat tertentu, model
disetel dengan reinforcement learning (\cref{ch:rl}), dengan model QSAR sebagai fungsi ganjaran, atau
dilatih ulang pada molekul-molekul terbaik. VAE \citep{gomezbombarelli2018} memetakan molekul ke ruang
laten kontinu, sehingga optimasi bisa dilakukan dengan gradien di ruang itu.

Ada beberapa jebakan. Model yang dioptimasi terhadap satu skor cenderung menemukan celah dalam skor itu,
menghasilkan molekul yang secara angka sempurna tetapi aneh secara kimia. Mode collapse membuat model
terus menghasilkan variasi kecil dari satu molekul. Untuk mengevaluasi model generatif, kita memeriksa
validitas, keunikan, kebaruan, dan keragaman molekul yang dihasilkan, serta seberapa mirip distribusinya
dengan molekul nyata. Fréchet ChemNet Distance \citep{preuer2018}, dari Linz juga, mengukur hal terakhir
dengan membandingkan distribusi aktivasi sebuah jaringan prediksi bioaktivitas untuk kedua himpunan
molekul. Kalau kedua distribusi dianggap Gauss, jarak Fréchet di antaranya punya bentuk tertutup,
\begin{equation}
d^2=\|\bm\mu_1-\bm\mu_2\|^2+\mathrm{Tr}\big(\bm\Sigma_1+\bm\Sigma_2-2(\bm\Sigma_1\bm\Sigma_2)^{1/2}\big).
\end{equation}
Dalam satu dimensi, rumus ini menyederhana menjadi $(\mu_1-\mu_2)^2+(\sigma_1-\sigma_2)^2$, karena
$\sigma_1^2+\sigma_2^2-2\sigma_1\sigma_2=(\sigma_1-\sigma_2)^2$: selisih rata-rata ditambah selisih
lebar sebaran. Kuliah juga mengingatkan bahwa metrik ini menangkap banyak bias dari jaringan yang
dipakainya, sehingga angka FCD yang baik belum menjamin molekul yang baik.

\section{Data medis}

Satu bagian kuliah membahas data klinis. Catatan kesehatan elektronik berisi data tabular dengan
campuran variabel kontinu, kategorikal, dan waktu, dengan banyak nilai yang hilang. Nilai yang hilang
jarang hilang secara acak: tes yang tidak dilakukan sering berarti dokter tidak menganggapnya perlu, dan
fakta itu sendiri membawa informasi. Kuliah membahas contoh model yang belajar sinyal yang salah, seperti
penanda rumah sakit atau jenis alat pemindai, alih-alih sinyal penyakit.

Untuk data tabular, model berbasis pohon masih sulit dikalahkan, meski arsitektur khusus terus
dikembangkan. TabPFN \citep{hollmann2025} mengambil jalan yang tidak biasa: sebuah transformer dilatih
sekali pada jutaan masalah tabular sintetis, lalu untuk data baru ia memprediksi langsung dalam satu
lintasan maju, dengan data latih sebagai konteks, tanpa pelatihan ulang. Data medis juga tersebar di
banyak rumah sakit yang tidak boleh saling berbagi data pasien. \istilah{Federated learning} melatih model
bersama dengan hanya bertukar pembaruan parameter, bukan data mentahnya.

\section{Merencanakan sintesis}

Molekul yang dirancang komputer tidak berguna kalau tidak bisa dibuat. \istilah{Retrosynthesis} bekerja
mundur: dari molekul target, cari reaksi yang bisa menghasilkannya dari bahan yang lebih sederhana, lalu
ulangi sampai tiba di bahan yang bisa dibeli. Satu langkah mundur bisa dirumuskan sebagai klasifikasi
atas aturan reaksi yang diekstrak dari basis data, sebagai terjemahan dari SMILES produk ke SMILES
reaktan, atau sebagai pencarian aturan yang paling cocok. MHNreact \citep{seidl2022} memakai jaringan
Hopfield modern untuk mencocokkan molekul dengan templat reaksi, sehingga templat yang jarang muncul pun
bisa dipilih. Untuk rencana banyak langkah, masalahnya mirip permainan: setiap langkah membuka banyak
kemungkinan, dan kita mencari jalur yang berakhir di bahan yang tersedia. \citet{segler2018} memakai
Monte Carlo tree search yang dipandu jaringan saraf, gagasan yang sama dengan AlphaGo (\cref{ch:klasik}).

\section{Benang merah}

Bab ini memperlihatkan bahwa hampir semua alat dari bagian-bagian sebelumnya menemukan tempatnya di
kimia dan kedokteran. Message passing dan ekuivariansi untuk molekul, LSTM dan transformer untuk SMILES
dan urutan, contrastive learning untuk menghubungkan modalitas, reinforcement learning dan pencarian
pohon untuk desain dan sintesis, jaringan Hopfield untuk belajar dari sedikit contoh. Yang membedakan
bidang ini adalah datanya: sedikit, berisik, tidak seimbang, dan penuh struktur tersembunyi yang bisa
membocorkan jawaban. Di sini, cara mengevaluasi sering lebih menentukan daripada arsitektur model.

\sumberbab{Bab ini mengikuti kuliah AI in Life Sciences: penemuan obat, representasi molekul, aktivitas
biologis dan ADME, QSAR dan virtual screening, evaluasi dan pembagian data, geometric deep learning untuk
molekul, model multimodal dan few-shot, model generatif, data medis, dan perencanaan sintesis.}
